\documentclass[11pt,a4paper]{article}
\usepackage[margin=25mm]{geometry}
\usepackage{amsmath,amssymb,amsthm,booktabs,array}
\pdfmapfile{+cm.map}
\pdfmapfile{+symbols.map}
\usepackage[colorlinks=true,linkcolor=blue,urlcolor=blue,citecolor=blue]{hyperref}
\usepackage[expansion=false]{microtype}
\hypersetup{pdftitle={Sharp all spin dimension bounds for finite Veneziano products},pdfauthor={Yicheng Pan},pdfsubject={Public AI-assisted preprint}}
\setlength{\emergencystretch}{2em}
\renewcommand{\labelitemi}{$\bullet$}
\newtheorem{theorem}{Theorem}[section]
\newtheorem{lemma}[theorem]{Lemma}
\newtheorem{proposition}[theorem]{Proposition}
\newtheorem{corollary}[theorem]{Corollary}
\theoremstyle{remark}\newtheorem{remark}[theorem]{Remark}
\newcommand{\Res}{\operatorname{Res}}
\newcommand{\E}{\mathbb E}
\newcommand{\dd}{\mathrm d}
\newcommand{\NN}{\mathbb N}
\newcommand{\cset}{\mathcal N}
\title{Sharp all spin dimension bounds for finite Veneziano products\\[5pt]\large A uniform gap and exact positivity certificates}
\author{Yicheng Pan}
\date{3 October 2026\\[3pt]\small Version 2 public AI-assisted preprint}
\begin{document}
\maketitle
\begin{abstract}
We determine the tree-level four-point partial-wave positivity range of every finite, equidistant Veneziano product. For each integer $n\geq3$, positivity at all mass levels and all spins holds exactly for real dimensions $3<D\leq d_n$, where $d_n$ is the unique zero of the level-three scalar projection. The thresholds decrease strictly to ten and satisfy $d_n=10+72/(11n)+13032/(1331n^2)+O(n^{-3})$. The uniform step is a gap theorem at $D=51/5$: all central-factorial Veneziano residues are positive except one low-level scalar. An established Bessel and first-exit method reduces this statement to a fixed integer polynomial with 182 positive coefficients. It settles all products with $n\geq35$; exact rational certificates cover the finite base $3\leq n\leq34$. The result provides a proof of a published numerical observation. The amplitude family, threshold values, and general proof strategies are prior work; originality of the specific uniform certificate and application remains provisional. No higher-point, loop-unitary, local, or physical string completion is established.
\end{abstract}
\paragraph{AI disclosure.} OpenAI AI assistance was used substantially for research, derivation, writing, and code. Separately executed AI-generated checks reconstructed the algebra and exact certificates. These are not human expert review or peer review. The named author's personal review is not asserted. This public AI-assisted preprint asserts no affiliation or contact information.

\section{Question and relation to existing work}
The finite planar products below are established infinite-spin-tower amplitudes. Shao and Vichi construct them in Ref.~\cite{Shao}. Section IV.3 of arXiv:2607.27300v2 gives numerical critical dimensions, a hypergeometric level-three scalar formula, and the limit ten. Its Table 1 suggests that this scalar is the first partial wave to become negative for every truncation. We prove that criterion for every integer number of massive poles, while retaining the precise scope of four-point tree-level partial-wave positivity.

The proof separates an infinite parameter range from a finite base. A uniform positivity gap for the underlying polynomial residues settles all $n\geq35$. Exact rational certificates cover
\begin{equation}\label{eq:set}
\cset=\{3,4,\ldots,34\}.
\end{equation}
The reduction to this finite set is proved analytically; it is not an inference from a larger numerical scan.

The tools and strategies have close precedents. Positive power-series coefficients are used for Gegenbauer positivity in Ref.~\cite[Section 4.6.1]{Geiser}. Type-I Bessel representations appear in Ref.~\cite[Appendix B]{AHHM}; that work also derives the relevant low-level scalar polynomial and discusses asymptotic positivity without a rigorous uniform cutoff. Chen and Yin's bosonic proof~\cite{Chen} supplies the central-factorial and finite-width first-exit method adapted here. Its kinematic parameter is $\kappa=c+2$, whereas our proof requires $\kappa=c$.

Mansfield~\cite[Section 3.2]{Mansfield} already uses a supercritical gap with one exceptional low level for a different massless hypergeometric amplitude. Numerical low-spin dominance in another family appears in Ref.~\cite[Section 4.1.2]{Wang}, and above-ten positivity occurs in Ref.~\cite{Cheung}. Neither these ideas nor their general strategies are claimed as new. The candidate contribution is the particular uniform $D=51/5$ inequality and fixed certificate, its application to all finite equidistant products, and the accompanying scalar analysis.

A bounded search of current versions and closely related papers located no matching proof of these exact statements. This supports specialist assessment, not a worldwide priority guarantee. This version extends an earlier finite-set research draft, whose scope was limited to $n=3,\ldots,16,32$.

\section{Amplitude and positivity conventions}
Normalize the first massive mass-squared to one and define
\begin{equation}\label{eq:amplitude}
 A_n(s,t)=-\frac{s+t}{st}\prod_{r=1}^{n}
 \frac{r(r-s-t)}{(r-s)(r-t)}.
\end{equation}
There are massive poles at $s=1,\ldots,n$. To fix the overall sign unambiguously, define the spectral residue by
\begin{equation}\label{eq:residue}
 R_{n,k}(t):=-\Res_{s=k}A_n(s,t)
 =\frac{(t+1)_k}{k!}\prod_{i=0}^{k-1}\frac{n-i}{n-i-t},
 \qquad 1\leq k\leq n.
\end{equation}
Here $(a)_m$ denotes a rising factorial. The minus sign in \eqref{eq:residue} is relative to the mathematical complex-analytic residue. At $s=0$ the spectral residue is the positive constant one.

At level $k$, use the massless angular variable
\begin{equation}\label{eq:angle}
 x=1+\frac{2t}{k},\qquad -1\leq x\leq1,
 \qquad \lambda=\frac{D-3}{2}>0.
\end{equation}
Tree-level partial-wave positivity means nonnegative coefficients in the expansion of $R_{n,k}$ in $C_j^\lambda(x)$. Replacing these polynomials by their normalization at $x=1$ changes coefficients only by positive factors. Every residue obeys $R_{n,k}(1)=1$.

Let $\E_D$ denote expectation with respect to the probability measure
\begin{equation}\label{eq:beta}
 \dd\nu_D(x)=\frac{(1-x^2)^{(D-4)/2}\dd x}
 {B(1/2,(D-2)/2)},\qquad
 M_m(D):=\E_D[X^{2m}]=\frac{(1/2)_m}{((D-1)/2)_m}.
\end{equation}
The normalized scalar projection is $F_n(D)=\E_D[R_{n,3}(X)]$. Its sign is the sign of the spin-zero coefficient. All results below refer to real $D>3$ in this precisely defined analytic continuation.

\begin{theorem}\label{thm:global}
For every integer $n\geq3$ and real $D>3$, all massive and massless spectral residues of \eqref{eq:amplitude} have nonnegative Gegenbauer coefficients at every integer spin $j\geq0$ if and only if $D\leq d_n$. Here $d_n$ is the unique level-three scalar zero specified in Theorem~\ref{thm:main}. The endpoint is included: nonnegativity allows its scalar coefficient to vanish.
\end{theorem}
The proof is completed in Section~\ref{sec:globalclosure}. It combines the scalar theorem, a uniform polynomial-residue gap, and an exact finite base.


\section{The uniform level three theorem}
\begin{theorem}\label{thm:main}
For each integer $n\geq3$, the function $F_n(D)$ has a unique zero $d_n\in(10,14)$. It is positive for $3<D<d_n$ and negative for $D>d_n$. All nonzero-spin level-three coefficients are strictly positive for $3<D\leq14$. Thus level three has nonnegative coefficients at every spin if and only if $3<D\leq d_n$.

The sequence $d_n$ is strictly decreasing and tends to ten. Moreover,
\begin{equation}\label{eq:asymptotic}
 d_n=10+\frac{72}{11n}+\frac{13032}{1331n^2}+O(n^{-3}).
\end{equation}
For $n=3$, the same interval $3<D\leq d_3$ is necessary and sufficient for positivity of the entire amplitude, including every level and spin.
\end{theorem}

\subsection{An infinite positive Taylor tail}
Set
\begin{align}
 q_i&=\frac{3}{2n-2i+3},\quad i=0,1,2,
 &K_n&=\frac1{16}\prod_{i=0}^{2}\frac{n-i}{n-i+3/2}>0,\label{eq:q}\\
 B_n(x)&=(1+x)\prod_{i=0}^{2}(1-q_i x)^{-1}
 =\sum_{m\geq0}b_m x^m.
\end{align}
Equation~\eqref{eq:residue} becomes
\begin{equation}\label{eq:R3}
 R_{n,3}(x)=K_n(9x^2-1)B_n(x)=\sum_{m\geq0}a_m x^m.
\end{equation}
All series here converge absolutely and uniformly on $[-1,1]$: their radius is at least $5/3$.

\begin{lemma}\label{lem:tail}
The coefficients in \eqref{eq:R3} satisfy $a_0=-K_n$, $a_1=-K_n b_1$, and $a_m>0$ for every $m\geq2$.
\end{lemma}
\begin{proof}
Let $h_m$ be the complete homogeneous symmetric polynomial of degree $m$ in the positive variables $q_0,q_1,q_2$, with $h_0=1$. Then $b_0=1$ and $b_m=h_m+h_{m-1}$ for $m\geq1$. Writing $S=q_0+q_1+q_2$, we have
\begin{equation}\label{eq:S}
 S\leq\frac13+\frac37+\frac35=\frac{143}{105},\qquad S^2<4.
\end{equation}
Multiplication of $h_{m-1}$ by $S$ counts every degree-$m$ monomial at least once, so $h_m\leq S h_{m-1}$. It follows that $b_m\leq S b_{m-1}$ for $m\geq2$. In particular,
\begin{equation}
 b_2\leq S^2+S<9,\qquad
 b_m\leq S^2 b_{m-2}<9b_{m-2}\quad(m\geq3).
\end{equation}
The claim follows from $a_m=K_n(9b_{m-2}-b_m)$ for $m\geq2$.
\end{proof}

Every monomial has a nonnegative Gegenbauer expansion for $\lambda>0$. One explicit identity is \cite[Eq.~18.18.17]{DLMF}
\begin{equation}\label{eq:monomial}
 x^m=\frac{m!}{2^m}\sum_{r=0}^{\lfloor m/2\rfloor}
 \frac{\lambda+m-2r}{\lambda}
 \frac{C_{m-2r}^{\lambda}(x)}{(\lambda+1)_{m-r}\,r!}.
\end{equation}
Consequently every spin $j\geq2$ in \eqref{eq:R3} is strictly positive. The negative terms have degree at most one and cannot contribute to these spins, while $a_j x^j$ gives a strictly positive contribution. Uniform convergence justifies projection. More explicitly, endpoint normalization of \eqref{eq:monomial} bounds the total absolute Gegenbauer weight by $\sum_m|a_m|<\infty$.

For spin one, the sign is that of $\E_D[X R_{n,3}(X)]$. Keeping only the negative linear term and the positive cubic term gives
\begin{align}\label{eq:spinone}
 \E_D[X R_{n,3}(X)]
 &\geq a_1M_1(D)+a_3M_2(D)\notag\\
 &\geq\frac{K_n b_1}{D-1}
 \left[-1+\frac{3(9-S^2)}{D+1}\right]>0
 \qquad(3<D\leq14).
\end{align}
The last inequality follows from $S^2<4$. This proves the infinite-spin part of Theorem~\ref{thm:main}.

\subsection{The scalar root and its monotonicity}
Lemma~\ref{lem:tail} gives
\begin{equation}\label{eq:Fdecreasing}
 F_n(D)=-K_n+\sum_{m\geq1}a_{2m}M_m(D).
\end{equation}
Every $M_m(D)$ for $m\geq1$ decreases strictly with $D>3$. Absolute summability of the coefficients gives continuity, strict decrease of $F_n$, and the limit $F_n(D)\to-K_n$ as $D\to\infty$.

A second organization of the same expectation is
\begin{equation}\label{eq:Falternate}
 \frac{F_n(D)}{K_n}
 =\sum_{r\geq0}b_{2r}M_r(D)
 \frac{16r+10-D}{D+2r-1}.
\end{equation}
At $D=10$, the $r=0$ term vanishes and every other term is positive. Thus $F_n(10)>0$. Section~\ref{sec:scalarcertificate} gives an exact rational certificate of $F_3(14)<0$.

For $D=14$, the $r=0$ term of \eqref{eq:Falternate} is negative and independent of $n$, while all $r\geq1$ terms are positive. Each $b_{2r}(n)$ with $r\geq1$ strictly decreases as $n$ increases: it is a positive, nonconstant polynomial in the $q_i$, and all three $q_i$ decrease. Therefore $F_n(14)/K_n\leq F_3(14)/K_3<0$. Strict decrease in $D$ proves existence and uniqueness of $d_n\in(10,14)$.

For $n'>n$, evaluate \eqref{eq:Falternate} at $D=d_n$. Only its $r=0$ term is negative; all other factors are positive. Decreasing the $b_{2r}$ strictly decreases the sum, so $F_{n'}(d_n)<0$ and $d_{n'}<d_n$. Finally, $B_n(x)\to1+x$ uniformly, and
\begin{equation}
 \frac{F_n(D)}{K_n}\longrightarrow\frac{10-D}{D-1}.
\end{equation}
For every fixed $D>10$ the limit is negative. Together with $d_n>10$ and monotonicity, this proves $d_n\downarrow10$.

\subsection{The complete three mass product}
The two lower massive residues are
\begin{align}\label{eq:lower}
 R_{n,1}(x)&=\frac{n(1+x)}{2n+1-x},\\
 R_{n,2}(x)&=\frac{n(n-1)x(x+1)}{2(n+1-x)(n-x)}.
\end{align}
Both have nonnegative convergent Taylor coefficients and therefore nonnegative partial waves in every $D>3$. The massless spectral residue is one. At $n=3$ these, together with level three, exhaust the spectrum. At $D=d_3$ the zero scalar coefficient is allowed, so the endpoint is included. The scalar is negative at every $D>d_3$.

Among physical integer dimensions $D\geq4$, this particular four-point positivity test permits $D=4,\ldots,13$. This statement neither constructs a thirteen-dimensional string theory nor changes the critical dimension of an established string theory.

\section{Rate of approach to ten}
Put $z=1/n$ and $t=3(x-1)/2$. With
\begin{equation}
 V(x)=\frac{(9x^2-1)(x+1)}{16},
\end{equation}
the residue admits the uniform expansion
\begin{align}
 R_{n,3}(x)
 &=V(x)\prod_{i=0}^{2}\frac{1-iz}{1-(i+t)z}\notag\\
 &=V(x)\left[1+3tz+(6t^2+3t)z^2+O(z^3)\right].
\end{align}
Taking beta moments gives $F_n(D)=f_0(D)+z f_1(D)+z^2f_2(D)+O(z^3)$, where
\begin{align}
 f_0(D)&=-\frac{D-10}{16(D-1)},\\
 f_1(D)&=\frac{9(D-8)(D-2)}{32(D-1)(D+1)},\\
 f_2(D)&=-\frac{9(D-8)(D-2)}{16(D-1)(D+1)}.
\end{align}
Substitution of $D=10+c_1z+c_2z^2+O(z^3)$ and matching powers of $z$ gives $c_1=72/11$ and $c_2=13032/1331$.

The remainder is justified analytically. For $t\in[-3,0]$ and $i=0,1,2$, one has $|i+t|\leq3$, so the finite correction product is uniformly analytic on every closed disk $|z|\leq\rho<1/3$. The normalized beta integral is jointly analytic near $(z,D)=(0,10)$, including complex $D$ in a neighborhood where endpoint weights remain integrable. Since $f_0'(10)=-1/144\neq0$, the analytic implicit-function theorem supplies a local root expansion. Uniqueness identifies this root with $d_n$ for all sufficiently large integers $n$, establishing \eqref{eq:asymptotic}.

By Theorem~\ref{thm:global}, this is also the asymptotic expansion of the full four-point positivity threshold for every finite product in the family.

\section{An exact scalar certificate}\label{sec:scalarcertificate}
At $n=3$, partial fractions give
\begin{equation}\label{eq:R33}
 R_{3,3}(x)=-1+\frac{24}{5-3x}-\frac{120}{7-3x}+\frac{40}{3-x}.
\end{equation}
Thus $a_0=-1/105$ and, for $m\geq1$,
\begin{equation}\label{eq:am}
 a_m=\frac{24}{5}\left(\frac35\right)^m
 -\frac{120}{7}\left(\frac37\right)^m
 +\frac{40}{3}\left(\frac13\right)^m.
\end{equation}
For rational $D>3$ and integer $J\geq0$, define
\begin{align}
 L_J(D)&=\sum_{j=0}^{J}a_{2j}M_j(D),\\
 T_J&=\frac{24}{5}\frac{(9/25)^{J+1}}{1-9/25}
       +\frac{40}{3}\frac{(1/9)^{J+1}}{1-1/9}.
\end{align}
All quantities are rational. Lemma~\ref{lem:tail}, the bound $M_j(D)\leq1$, and omission of the negative term in \eqref{eq:am} yield
\begin{equation}\label{eq:interval}
 L_J(D)<F_3(D)<L_J(D)+T_J.
\end{equation}
Using $J=100$ gives $F_3(10)>0.0055688$ and $F_3(14)<-0.00089865$. More precisely, the same exact-arithmetic test proves
\begin{align}\label{eq:bracket}
 13.15983426171950364948847935837366808&<d_3,\\
 d_3&<13.15983426171950364948847935837366809.\notag
\end{align}
The lower endpoint scalar is greater than $9.11210\times10^{-39}$; the upper endpoint scalar is less than $-2.57466\times10^{-39}$. The enclosure width for either scalar is below $1.16\times10^{-44}$. These signs use exact rational arithmetic, rather than a floating-point root finder. The decimal precision is a reproducibility detail; the numerical threshold was already reported in Ref.~\cite{Shao}.

For a separate numerical cross-check, \eqref{eq:R33} gives the normalized scalar
\begin{align}
 F_3(D)={}&-1+\frac{24}{5}\,{}_2F_1\left(1,\tfrac12;\tfrac{D-1}{2};\tfrac9{25}\right)\notag\\
 &-\frac{120}{7}\,{}_2F_1\left(1,\tfrac12;\tfrac{D-1}{2};\tfrac9{49}\right)
 +\frac{40}{3}\,{}_2F_1\left(1,\tfrac12;\tfrac{D-1}{2};\tfrac19\right).
\end{align}
A 90-digit evaluation agrees with direct normalized beta-weight quadrature. These numerical checks support the certificate but do not replace it.


\section{A uniform gap for polynomial residues}\label{sec:gap}
For nonnegative integers $m$ define
\begin{equation}\label{eq:central}
 Q_m(y)=\prod_{a=-(m-1)/2}^{(m-1)/2}(y-a),\qquad
 W_m(x)=\frac1{m!}Q_m\!\left(\frac{(m+1)x}{2}\right),
\end{equation}
where the product advances in unit steps and $Q_0=1$. The polynomial factor of \eqref{eq:residue} satisfies
\begin{equation}\label{eq:VW}
 V_k(x):=\frac{(t+1)_k}{k!}
 =\frac{1+x}{2}W_{k-1}(x),\qquad t=\frac{k}{2}(x-1).
\end{equation}

\begin{theorem}\label{thm:gap}
At $D_*=51/5$, every parity-allowed Gegenbauer coefficient of $W_m$ is strictly positive except $(m,j)=(2,0)$, which is negative. Coefficients with $j>m$ or $m-j$ odd vanish. In particular, every $W_m$ with $m\neq2$ has nonnegative coefficients for all real $3<D\leq51/5$.
\end{theorem}
The integers $m,j\geq0$ are essential to the exception statement. Dimension descent follows from \eqref{eq:connection}; it remains to prove the endpoint $D_*$.

\subsection{Central factorials and the Bessel transform}
The established central-factorial identity~\cite{Chen} is
\begin{equation}\label{eq:centralGF}
 \sum_{m\geq0}\frac{Q_m(y)}{m!}u^m
 =\frac{\exp(2y\,\operatorname{arsinh}(u/2))}{\sqrt{1+u^2/4}}.
\end{equation}
For completeness, set $u=2\sinh z$ and then $w=e^{2z}$ in the coefficient contour. The coefficient of $u^m$ becomes
\begin{equation}
 \Res_{w=1}\frac{w^{y+(m-1)/2}}{(w-1)^{m+1}}
 =\binom{y+(m-1)/2}{m}=\frac{Q_m(y)}{m!}.
\end{equation}
For matching parity write
\begin{equation}\label{eq:Besselparams}
 \ell=\frac{m-j}{2},\quad c=m+1=j+2\ell+1,\quad
 b=j+\frac{D-1}{2},\quad \lambda=\frac{D-3}{2}.
\end{equation}
Insert the Gegenbauer plane-wave expansion into \eqref{eq:centralGF} and again change $u=2\sinh z$. Its Jacobian cancels the reciprocal $\cosh z$ factor. The coefficient $B_{m,j}(D)$ of $C_j^\lambda$ in $W_m$ is
\begin{align}\label{eq:BesselH}
 H(z)&=\left(\frac{z}{\sinh z}\right)^c
 {}_0F_1\!\left(;b;\frac{c^2z^2}{4}\right)
 =\sum_{p\geq0}h_pz^{2p},\\
 B_{m,j}(D)&=\frac{c^j}{2^{m+j}(\lambda)_j}\,h_\ell.\label{eq:prefactor}
\end{align}
The prefactor is positive. This is the same type of transform used in Refs.~\cite{AHHM,Chen}, with the present parameter $\kappa=c$ fixed by \eqref{eq:central}.

\subsection{The recurrence and positive bulk}
The hypergeometric factor $\Phi={}_0F_1(;b;c^2z^2/4)$ obeys
\begin{equation}
 \Phi''+\frac{2b-1}{z}\Phi'-c^2\Phi=0.
\end{equation}
Substituting $\Phi=(\sinh z/z)^cH$ and multiplying by $z^2\sinh^2z$ gives
\begin{equation}\label{eq:conjugated}
\begin{split}
 z^2\sinh^2z\,H^{\prime\prime}&+2cz^2\sinh z\cosh z\,H^{\prime}\\
 &+(D-4-4\ell)z\sinh^2z\,H^{\prime}=N(z)H,\\
 N(z)&=c^2z^2\sinh^2z-c^2z^2\cosh^2z+cz^2\\
 &\quad+c(2c-2b+1)z\sinh z\cosh z\\
 &\quad+c(2b-c-2)\sinh^2z.
\end{split}
\end{equation}
Define $\alpha_r=2^{2r-1}/(2r)!$. For $p\geq1$, coefficient extraction yields
\begin{align}\label{eq:uniformrec}
 \Delta_ph_p&=\sum_{r=2}^{p+1}\alpha_rM_{p,r}h_{p+1-r},
 &\Delta_p&=2p(2j+2p+D-3),\\
 M_{p,r}&=c[(4\ell+4-D)r+j-2\ell+D-4]\notag\\
 &\quad-2q[2q-1+2cr+D-4-4\ell],
 &q&=p+1-r.\label{eq:uniformM}
\end{align}
The target parameters $c,b,\ell$ remain fixed as the recurrence index $p$ decreases. Resetting them at each recurrence step would be invalid.

Throughout the remaining bulk, first-exit, and shallow-trajectory calculations in this section, set $D=D_*=51/5$. For $p\leq\ell-1$ set $\ell-p=1+X$, $r=2+Y$, and $p=r-1+Q$, with $X,Y,Q\geq0$. Then \eqref{eq:uniformM} becomes
\begin{align}\label{eq:bulkpoly}
 M_{p,r}={}&8QXY+16QX+8QY^2+\tfrac{98}{5}QY+\tfrac{36}{5}Q\notag\\
 &+8X^2Y+12X^2+16XY^2+4XYj+\tfrac{238}{5}XY\notag\\
 &+8Xj+\tfrac{208}{5}X+8Y^3+4Y^2j+\tfrac{178}{5}Y^2\notag\\
 &+\tfrac{49}{5}Yj+\tfrac{253}{5}Y+j^2+\tfrac{54}{5}j+29.
\end{align}
All twenty coefficients are positive. Since $h_0=1$ and $\Delta_p>0$, induction proves $h_p>0$ for $0\leq p\leq\ell-1$. Some edges at the final target $p=\ell$ are negative; that target requires a separate identity.

\subsection{A fixed three step boundary certificate}
Following the first-exit construction of Ref.~\cite{Chen}, define
\begin{align}
 C_{p,r}&=\frac{\alpha_rM_{p,r}}{\Delta_p},&
 F_0&=1,&
 F_d&=\sum_{u=0}^{d-1}F_uC_{\ell-u,d-u+1}\quad(d=1,2),\\
 E_R&=\sum_{u=0}^{2}F_uC_{\ell-u,R-u+1}
 \quad(3\leq R\leq\ell).
\end{align}
For $\ell\geq3$, recursively expand only at depths zero, one and two below $h_\ell$, stopping on first exit to depth $R\geq3$. The finite acyclic recurrence gives exactly
\begin{equation}\label{eq:firstexit}
 h_\ell=\sum_{R=3}^{\ell}E_Rh_{\ell-R}.
\end{equation}
All landing coefficients belong to the positive bulk.

The exit signs reduce to a fixed integer polynomial. To define it without a long coefficient listing, put $D_u=\Delta_{\ell-u}$, $T_0=1$, and
\begin{equation}
 T_d=\sum_{u=0}^{d-1}T_u\alpha_{d-u+1}M_{\ell-u,d-u+1}
       \prod_{v=u+1}^{d-1}D_v\quad(d=1,2).
\end{equation}
Thus $T_d=F_d\prod_{v=0}^{d-1}D_v$. Using
$\alpha_{R-u+1}/\alpha_{R+1}=4^{-u}(2R+2)^{\underline{2u}}$, define
\begin{equation}\label{eq:Pconstruction}
 \mathcal P(j,X,Y)=
 \left.
 22500\sum_{u=0}^{2}
 T_u\,4^{-u}(2R+2)^{\underline{2u}}
 M_{\ell-u,R-u+1}\prod_{v=u+1}^{2}D_v
 \right|_{\substack{R=X+3\\\ell=X+Y+3}} .
\end{equation}
The underlined exponent denotes a falling factorial. Exact expansion gives
\begin{align}\label{eq:Pcertificate}
 &\mathcal P\in\mathbb Z[j,X,Y],\qquad
 |\operatorname{supp}\mathcal P|=182,\notag\\
 &\deg\mathcal P=11,\qquad
 (\deg_j,\deg_X,\deg_Y)=(6,10,5),\notag\\
 &\min\{\text{nonzero coefficients}\}=2500,\qquad
 \mathcal P(0,0,0)=3381984900,
\end{align}
and every nonzero coefficient is positive. The full coefficient list and a short exact reconstruction program accompany the manuscript. This finite identity implies positivity on the entire nonnegative orthant, not merely at sampled levels or spins. In particular,
\begin{equation}\label{eq:Epositive}
 E_R=\frac{\alpha_{R+1}\mathcal P(j,R-3,\ell-R)}
 {22500\,\Delta_\ell\Delta_{\ell-1}\Delta_{\ell-2}}>0.
\end{equation}
Equations \eqref{eq:firstexit} and \eqref{eq:Epositive} prove $h_\ell>0$ for all $\ell\geq3$.

\subsection{Shallow trajectories and the single exception}
For the separate target cases $\ell=0,1,2$, set $c=j+2\ell+1$ in each case. Direct expansion of \eqref{eq:BesselH} at the respective target gives
\begin{align}
 h_0&=1,\qquad
 h_1=\frac{(j+3)(5j-1)}{12(5j+23)},\\
 h_2&=\frac{(j+5)(125j^3+1275j^2+3295j+177)}
 {1440(5j+23)(5j+28)}.
\end{align}
For integer $j\geq0$, $h_1$ is negative only at $j=0$, corresponding to $(m,j)=(2,0)$, while $h_2$ is positive. Together with the boundary certificate and \eqref{eq:prefactor}, this proves Theorem~\ref{thm:gap}.

The convenient value $51/5$ is not asserted to be the largest possible gap. For comparison, the level-five scalar has its lower zero at $(125-\sqrt{1081})/9=10.235715\ldots$, already encoded in Ref.~\cite[Table 1]{AHHM} up to a positive residue normalization. The simpler proposed gap $D=11$ is false: in the normalization \eqref{eq:VW}, that scalar equals $-1/2048$.

\section{The exact finite base}\label{sec:base}
\begin{theorem}\label{thm:finite}
For every $n\in\cset$ in \eqref{eq:set}, the amplitude \eqref{eq:amplitude} has nonnegative coefficients at all mass levels and all spins if and only if $3<D\leq d_n$.
\end{theorem}
The proof consists of the analytic reductions below and finitely many exact rational sign certificates, reproduced by the accompanying program. Theorem~\ref{thm:main} treats $k=3$; the massless level is automatic.

\subsection{Partial fractions and an infinite tail}
For $1\leq k\leq n$, put $a=n-k+1$ and define, for $r=a,\ldots,n$,
\begin{align}
 z_r&=1+\frac{2r}{k},\qquad q_r=z_r^{-1},\\
 A_r&=\frac2k\binom{r+k}{k}\frac{n!}{(n-k)!}
       \frac{(-1)^{r-a}}{(r-a)!(n-r)!}.
\end{align}
Taking the simple-pole residues and the limit at infinity gives the algebraic identity
\begin{equation}\label{eq:pf}
 R_{n,k}(x)=(-1)^k\binom nk+\sum_{r=a}^{n}\frac{A_r}{z_r-x}.
\end{equation}
For $m\geq1$, its Taylor coefficient is $u_m=\sum_r A_r q_r^{m+1}$; at $m=0$ add the displayed constant.

The nearest pole has $A_a>0$ and $q_a>q_r$ for every $r>a$. Choose $M_0\geq1$ such that
\begin{equation}\label{eq:dominance}
 A_aq_a^{M_0+1}>\sum_{A_r<0}|A_r|q_r^{M_0+1}.
\end{equation}
Such an integer exists. Dividing by $q_a^{m+1}$ shows that the right-to-left ratio only decreases for $m\geq M_0$. Hence $u_m>0$ for every $m\geq M_0$.

Compute $u_0,\ldots,u_{M_0}$ exactly and let $L$ be the largest degree with $u_L<0$, taking $L=-1$ if there is none. Equation~\eqref{eq:monomial} then proves positivity for all spins $j>L$. This is the step controlling infinitely many spins.

\subsection{Finite rational lower bounds}
For $0\leq j\leq L$, the Gegenbauer Rodrigues formula and $j$ integrations by parts show that the partial-wave sign is the sign of
\begin{equation}\label{eq:G}
 G_j(D)=\sum_{\substack{m\geq j\\m-j\ \mathrm{even}}}
 u_m\binom mj\,M_{(m-j)/2}(D+2j).
\end{equation}
The omitted proportionality factor is positive. Equivalently, \eqref{eq:G} follows from the ratios of coefficients in \eqref{eq:monomial}. The rational Taylor series and each of its derivatives converge uniformly on $[-1,1]$, so termwise operations are permitted. In the integration-by-parts derivation, the weight still vanishes at both endpoints after at most $j-1$ derivatives for every $\lambda>0$.

Take $M\geq M_0$. All terms omitted from \eqref{eq:G} after degree $M$ are nonnegative. A strictly positive partial sum through $M$ is therefore a rigorous lower bound on the complete spin-$j$ projection. At rational $D$ the partial sum is rational. No extrapolation in spin and no floating-point sign decision is needed.

\subsection{Dimension descent and the certificate set}
For each $n\in\cset$, choose the rational cap $D_{\rm cap}$ in Table~\ref{tab:caps}. A level-three scalar enclosure, obtained from \eqref{eq:pf} and a geometric tail bound as in \eqref{eq:interval}, proves $F_n(D_{\rm cap})<0$. Thus $d_n<D_{\rm cap}$. All residues $k\neq3$ pass \eqref{eq:dominance} and the finite lower-bound tests \eqref{eq:G} at that cap.

Positivity persists under decreasing dimension. For $\mu\geq\lambda>0$, the connection formula \cite[Eq.~18.18.16]{DLMF} is
\begin{equation}\label{eq:connection}
 C_j^\mu(x)=\sum_{\ell=0}^{\lfloor j/2\rfloor}
 \frac{\lambda+j-2\ell}{\lambda}
 \frac{(\mu)_{j-\ell}}{(\lambda+1)_{j-\ell}}
 \frac{(\mu-\lambda)_\ell}{\ell!}
 C_{j-2\ell}^{\lambda}(x).
\end{equation}
Every coefficient is nonnegative. Set $\mu=(D_{\rm cap}-3)/2$ and $\lambda=(D-3)/2$. Absolute convergence follows from the analytic residue and its finite positive endpoint weight; alternatively use $\sum_m|u_m|<\infty$ and \eqref{eq:monomial}. The connection formula therefore transfers the positive expansion to every $3<D\leq D_{\rm cap}$.

All $k\neq3$ residues are consequently positive at $D\leq d_n$. Theorem~\ref{thm:main} supplies the remaining level and necessity above $d_n$, proving Theorem~\ref{thm:finite}.

\begin{table}[t]
\centering
\caption{Exact finite-base certificates. Each cap is $D_{\rm cap}=p_n/2^{17}>d_n$. The column $M$ is the largest finite Taylor projection degree used at $k\neq3$; infinitely many remaining spins are handled analytically. Full per-residue data are provided in the certificate file.}\label{tab:caps}
\begin{tabular}{r r r @{\qquad} r r r}\toprule
$n$&$p_n$&$M$&$n$&$p_n$&$M$\\\midrule
3 & 1724886 & 64 & 19 & 1359217 & 64 \\
4 & 1594958 & 64 & 20 & 1356641 & 64 \\
5 & 1526895 & 64 & 21 & 1354324 & 66 \\
6 & 1484886 & 64 & 22 & 1352228 & 70 \\
7 & 1456388 & 64 & 23 & 1350324 & 74 \\
8 & 1435807 & 64 & 24 & 1348586 & 78 \\
9 & 1420259 & 64 & 25 & 1346993 & 82 \\
10 & 1408110 & 64 & 26 & 1345529 & 86 \\
11 & 1398361 & 64 & 27 & 1344178 & 90 \\
12 & 1390368 & 64 & 28 & 1342927 & 188 \\
13 & 1383698 & 64 & 29 & 1341766 & 99 \\
14 & 1378051 & 64 & 30 & 1340686 & 103 \\
15 & 1373208 & 64 & 31 & 1339678 & 214 \\
16 & 1369010 & 64 & 32 & 1338736 & 222 \\
17 & 1365337 & 64 & 33 & 1337853 & 230 \\
18 & 1362096 & 64 & 34 & 1337024 & 240 \\
\bottomrule\end{tabular}
\end{table}


\section{Closing the entire family}\label{sec:globalclosure}
The remaining step transfers Theorem~\ref{thm:gap} to every sufficiently long finite product. Rewrite \eqref{eq:residue} as
\begin{align}\label{eq:correction}
 R_{n,k}(x)&=V_k(x)\prod_{i=0}^{k-1}
              \frac{z_i-1}{z_i-x},
 &z_i&=1+\frac{2(n-i)}{k}>1.
\end{align}
Every rational correction factor has a positive Taylor series. Multiplication by $x$ has nonnegative matrix elements in the Gegenbauer basis for $\lambda>0$, so multiplication by any such Taylor polynomial preserves the positive cone. The factor $(1+x)/2$ in \eqref{eq:VW} also preserves it.

For each fixed finite pair $(n,k)$, the correction series converges uniformly on $[-1,1]$. Apply the cone argument to Taylor partial sums and then project the uniform limit. No uniform lower bound on angular pole separation as $n,k\to\infty$ is assumed. Thus Theorem~\ref{thm:gap} proves that every $k\neq3$ residue is nonnegative at all spins for $3<D\leq51/5$, uniformly in the allowed integer indices.

An exact scalar enclosure of the form used in Section~\ref{sec:scalarcertificate} gives
\begin{equation}\label{eq:cutoff35}
 F_{35}(51/5)<-\frac3{100000}<0.
\end{equation}
By Theorem~\ref{thm:main}, $d_n\leq d_{35}<51/5$ for every $n\geq35$. Hence all $k\neq3$ residues are positive for $D\leq d_n$, and the level-three theorem supplies the remaining residue. The massless residue is one. The exact base in Section~\ref{sec:base} covers $3\leq n\leq34$. For every $n$, $D>d_n$ is excluded by the negative scalar at level three. This completes the proof of Theorem~\ref{thm:global}, including the endpoint.

\begin{corollary}\label{cor:integer}
For integer dimensions $D\geq4$, the complete positivity ranges are: no upper bound for $n=1,2$; $D\leq13$ for $n=3$; $D\leq12$ for $n=4$; $D\leq11$ for $5\leq n\leq7$; and $D\leq10$ for $n\geq8$.
\end{corollary}
For $n=1,2$, the Taylor-positive formulas \eqref{eq:lower} exhaust the massive levels. For $n\geq3$, monotonicity of $d_n$ and exact scalar signs at $(n,D)=(3,13)$, $(3,14)$, $(4,12)$, $(4,13)$, $(5,12)$, $(7,11)$, and $(8,11)$ prove the stated transitions. These transitions were already suggested by the numerical table of Ref.~\cite{Shao}; the content here is their all-$n$ proof for this specific four-point family.

\section{Reproducibility and verification}
The bundle contains editable sources, a Chinese summary, exact certificates, and executable checks. The principal programs are:
\begin{itemize}
\item \path{verify_level3_theorem.py} and \path{derive_level3.py}: scalar enclosures, residue signs, and the asymptotic expansion.
\item \path{verify_uniform_gap.py}: exact reconstruction of the twenty-term bulk, the 182-term boundary polynomial, and the shallow trajectories. Its coefficient list is \path{uniform_gap_certificate.json}.
\item \path{certify_uniform_base.py}: all 560 non-level-three residue certificates, 32 scalar caps, and \eqref{eq:cutoff35}. Its data file is \path{uniform_base_certificates.json}.
\item \path{check_integer_dimension_corollary.py}: the seven scalar sign enclosures used in Corollary~\ref{cor:integer}.
\end{itemize}
The finite-base certifier uses standard-library rational arithmetic. Other principal calculations use SymPy and mpmath. No key sign decision uses floating-point tolerances. The largest Taylor truncation in the base is $M=240$; it is a finite certificate size, not a spin cutoff.

Separately written AI-generated checks derived the conjugated differential equation, reconstructed all first-exit paths, and regenerated the full finite base from polynomial long division rather than importing the principal implementation. A further adversarial check rebuilt the 182 coefficients in a custom rational sparse-polynomial ring without a computer algebra system and tested selected difficult base residues. The latter check did not independently replay all 560 rows; the full replay was performed by the former check. All stated checks agreed.

These are mathematical and computational cross-checks, not human expert review, peer review, or formal proof-assistant verification. Direct polynomial and numerical sanity checks are supplementary; the infinite claims rest on the symbolic reductions and exact certificates. The named author should review and approve the complete argument and disclosures before any submission.

\section{Scope and remaining questions}
The theorem concerns a known family of four-point meromorphic amplitudes with infinitely many spins at each massive pole. It establishes their tree-level partial-wave positivity range. It does not establish a local theory, a positive state-space construction, consistent higher-point factorization, loop unitarity, or a physical ultraviolet completion. The parameter $d_n$ is a threshold of this amplitude family, not a new critical dimension of string theory.

The amplitude construction, numerical thresholds, scalar hypergeometric formulas, and limiting value ten are prior results. Positive-monomial arguments, Bessel representations, exceptional-low-level strategies, and the first-exit mechanism are also prior methods. The possible original content is the particular uniform gap certificate and its all-$n$ application, together with the specific scalar inequalities and rate. A bounded literature audit found no exact matching proof, but originality and significance require human specialist assessment.

Further mathematical work could seek a more transparent positive representation that removes the finite-base computations, or determine the largest possible uniform gap after excluding level three. Those refinements would concern this positivity problem; a higher-point or loop completion would require additional independent arguments. This is a public AI-assisted preprint; no journal submission, acceptance, or arXiv endorsement is claimed.

\begin{thebibliography}{9}
\small\setlength{\itemsep}{3pt}
\bibitem{Shao}
L.-Q. Shao and A. Vichi,
\emph{Analytic Boundaries of Infinite-Spin-Tower Amplitudes from Hidden Zero},
\href{https://arxiv.org/abs/2607.27300v2}{arXiv:2607.27300v2} (2026), revised 31 August 2026. See Sections IV.1--IV.3 and Table 1.
\bibitem{Geiser}
N. Geiser and L. W. Lindwasser,
\emph{Generalized Veneziano and Virasoro amplitudes},
JHEP \textbf{04} (2023) 031,
\href{https://arxiv.org/abs/2210.14920v3}{arXiv:2210.14920v3}. See Sections 4.6.1--4.6.2.
\bibitem{Cheung}
C. Cheung and G. N. Remmen,
\emph{Veneziano Variations: How Unique are String Amplitudes?},
JHEP \textbf{01} (2023) 122,
\href{https://arxiv.org/abs/2210.12163v2}{arXiv:2210.12163v2}.
\bibitem{Chen}
Q. Chen and Y. Yin,
\emph{Exact All-Level Positivity of the Bosonic Veneziano Amplitude},
\href{https://arxiv.org/abs/2608.25183v1}{arXiv:2608.25183v1} (2026).
\bibitem{AHHM}
N. Arkani-Hamed, L. Eberhardt, Y.-t. Huang and S. Mizera,
\emph{On Unitarity of Tree-Level String Amplitudes},
JHEP \textbf{02} (2022) 197,
\href{https://arxiv.org/abs/2201.11575}{arXiv:2201.11575}.
See Appendix B, Table 1 and Section 4.3.1.
\bibitem{Mansfield}
G. Mansfield, \emph{Positivity of the Veneziano Amplitude in Ten Dimensions},
\href{https://arxiv.org/abs/2502.20372v2}{arXiv:2502.20372v2}, revised 17 February 2026. See Section 3.2.
\bibitem{Wang}
B. Wang, \emph{Positivity of the Hypergeometric Coon Amplitude},
JHEP \textbf{04} (2024) 143,
\href{https://arxiv.org/abs/2403.00906v4}{arXiv:2403.00906v4}.
See Section 4.1.2.
\bibitem{DLMF}
NIST Digital Library of Mathematical Functions,
\href{https://dlmf.nist.gov/18.18}{Section 18.18}, equations 18.18.16 and 18.18.17, accessed 3 October 2026.
\end{thebibliography}
\end{document}
